This paper presents a software-defined, modular railway interlocking system constructed with C++, Qt, and PostgreSQL \cite{elms2018interlocking, blanchette2006qt4}. The system provides real-time control of signals, switching of point machines and validation of route occupation while utilizing a declarative rule engine that evaluates JSON-based safety logic \cite{collier2020interlocking, mccluskey2017automated}. Because of the architecture, which allows for tight integration of the GUI, database and backend services with a consistent and synchronous process, simulation and hardware-in-the-loop (HIL) deployment are easy to conduct \cite{lehner2018hardware, chen2019modeling}. This system is designed using a hybrid model that facilitates resilience by employing database triggers with adaptive polling mechanisms, allowing for fault-tolerant updates, and ensuring continual monitoring of events \cite{costa2016resilient, peralta2021reliable}. Safety is built within the application, in multiple levels of validation cascades, system freeze workarounds, and formal error reporting \cite{bonacchia2016validation, li2020safety}. User activity is managed through a role-based access model that ensures critical controls are safeguarded \cite{sandhu1996role, li2020secure}. Other features of the system include rollback functionality for analysis of previous states, real-time logging for forensic-type inspection of record and issue identification, and a simulation interface for testing and operator training \cite{fernandez2018log, bennett2019observability}. Performance assessments of the system show responses well under 50 ms, and allows for repeatable identification of failed behavior predicated on fault conditions. The software design ensures a degree of extensibility and a secure foundation for developing digital interlocking elements from one-off railroad applications to broad input from varied physical layouts \cite{mccarthy2019towards, haxthausen2023compositional}.
